% Generated draft: requires mathematical and source review.
\documentclass[11pt]{article}
\usepackage[margin=25mm]{geometry}
\usepackage{fontspec}
\setmainfont{Latin Modern Roman}
\usepackage{amsmath,amssymb,mathtools,unicode-math}
\setmathfont{Latin Modern Math}
\usepackage{enumitem,needspace,xurl,hyperref}
\setlength{\parindent}{0pt}
\setlength{\parskip}{5pt}
\setlength{\emergencystretch}{3em}
\newcommand{\N}{\mathbb{N}}
\newcommand{\Z}{\mathbb{Z}}
\newcommand{\Q}{\mathbb{Q}}
\newcommand{\R}{\mathbb{R}}
\newcommand{\C}{\mathbb{C}}
\newcommand{\F}{\mathbb{F}}
\newcommand{\A}{\mathbb{A}}
\newcommand{\PP}{\mathbb{P}}
\newcommand{\E}{\mathbb{E}}
\newcommand{\eps}{\varepsilon}
\newcommand{\dd}{\,\mathrm{d}}
\begin{document}
\section*{UnsolvedMath ID: 20000827}
\textbf{Source problem code:} catalogue entry
\subsection{Definitions and mathematical statement}
Write $H_{d,g}$ for the locally Cohen--Macaulay curve locus in the Hilbert scheme of $\PP^3$ with Hilbert polynomial $dt+1-g$. A curve is locally Cohen--Macaulay if it has no embedded or isolated zero-dimensional components. Its Rao module is $M_C=\bigoplus_m H^1(\PP^3,\mathcal I_C(m))$. An extremal curve has the maximal Rao function allowed for its fixed degree and arithmetic genus, in the standard extremal-curve convention.

\textbf{Problem 26.} (i) If $C$ is a divisor of bidegree $(3,7)$ on a smooth quadric $Q\simeq\PP^1\times\PP^1$, can it be connected to an extremal curve in $H_{10,12}$? Its degree is $3+7=10$ and genus $(3-1)(7-1)=12$. (ii) Given four skew ruling lines on a smooth quadric $Q_1$, find a flat family of quadrics degenerating to a double plane $Q_0=2H$ and a flat family $C_t\subset Q_t$ whose special curve $C_0$ is locally Cohen--Macaulay. The original remark states (ii)$\Rightarrow$(i). The OCR arrow and footer damage are repaired according to the supplied clean text.

\textbf{Supplied resolution.} The aid attributes both affirmative answers to Lella--Schlesinger: for every $d\geq3$, $d$ disjoint ruling lines on a smooth quadric admit a flat degeneration to an extremal locally Cohen--Macaulay curve in a double plane, and every effective divisor on a smooth quadric belongs to the connected component containing extremal curves. In particular the $(3,7)$ curves are connected to the extremal locus in $H_{10,12}$. ``Connected component'' is essential; this does not say all such curves lie in the same irreducible component.

\textbf{Reported characteristic-zero specialization.} In coordinates $[x:y:z:w]$, for ruling parameters $\{0,1,3,7\}$ the aid gives the quadric family
\[
 Q_t:\ x^2+t^3xw-t^5yz=0,
\]
with central double plane $x^2=0$. Its reported weight degeneration has central curve ideal
\[
 I_E=(x^2,xy,y^4,xG-y^3F),
\]
\[
 G=(w+z)(w+3z)(w+4z)(w+7z)(w+8z)(w+10z),\qquad
 F=z^3(2w+11z).
\]
This is the claimed extremal length-four-ruling-lines limit, of degree $4$ and genus $-3$. The full equations for $C_t$ and the ruling-parameter convention are not supplied, so the displayed endpoint and quadrics do not alone reconstruct the entire flat family.

\textbf{Reported liaison audit.} A proper direct complete-intersection link of type $(2,7)$ from this degree-four curve would have residual degree $10$ and genus $12$. The aid proves that its Rao-module length is $24$, whereas a degree-$10$, genus-$12$ extremal curve has Rao-module length $384$. Hence that direct residual is not itself extremal. The implication from the four-line degeneration to connectivity of $(3,7)$ divisors genuinely uses biliaison and double-plane connectedness rather than this naive single-link identification. The existence result, explicit specialization and numerical obstruction are supplied claims requiring source review.
\subsection{Short English statement}
The aid reports positive solutions to both the quadric-divisor connectivity and four-skew-lines degeneration questions. It gives an explicit double-plane limit with its ideal. A Rao-length calculation also prevents identifying a direct $(2,7)$ residual with the desired extremal curve, so connectedness must use the stated biliaison argument.
\subsection{Sources}
\begin{itemize}
\item Local UnsolvedMath registry, permanent ID $20000827$, \texttt{catalogue entry}; canonical \texttt{aim-arithmetic-geometry-notes.json}, \texttt{notes[72]}.
\item AIM workshop ``Components of Hilbert schemes,'' PDF Problem 26 and remark (ii)$\Rightarrow$(i): \url{https://aimath.org/WWN/hilbertschemes/hilbertschemes.pdf}. The extracted footer says recorded on the whiteboard by Izzet Coskun and edited by Li Li; it is an editorial attribution, not part of the mathematical question. Registry statement status \texttt{reconstructed\_unverified}.
\item Lella--Schlesinger degeneration and connectedness results, as named in the attached summary; no full title, theorem number or bibliographic reference supplied.
\item Research attempt 1, \texttt{full\_solution}, run status \texttt{solved}, internal rigor \texttt{proved}, low novelty confidence; registry \texttt{SOLVED-IN-LITERATURE}. Explicitly machine-generated, not peer reviewed, and to be independently verified before citation. Original literature field empty.
\end{itemize}
\end{document}
